\documentclass[11pt]{beamer}
\usetheme{Madrid}
\usepackage[T1]{fontenc}
\usepackage[utf8]{inputenc}
\usepackage{amsmath,amsfonts,amssymb}
\usepackage{graphicx}
\usepackage{xcolor}
\usepackage{tikz}
\usepackage{array}
\usepackage{longtable}
\usepackage{booktabs}
\usepackage{hyperref}

\title{EB2-NIW \& EB1A Profile Building}
\subtitle{High-Volume Research Strategy: 20-50 Papers | 1000-2000 Citations}
\author{Complete DIY Research Profile System}
\date{\today}

\setbeamercovered{transparent}
\setbeamertemplate{navigation symbols}{}
\setbeamerfont{itemize/enumerate body}{size=\scriptsize}
\setbeamerfont{itemize/enumerate subbody}{size=\tiny}
\setbeamerfont{block body}{size=\scriptsize}
\setbeamerfont{block title}{size=\normalsize}

\begin{document}
	
	% ============================================================
	% SLIDE 1: TITLE
	% ============================================================
	\begin{frame}
		\titlepage
		\centering
		\scriptsize \textit{20-50 Papers | 1000-2000 Citations | Government Impact Focus}
	\end{frame}
	
	% ============================================================
	% SLIDE 2: THE TARGET NUMBERS
	% ============================================================
	\begin{frame}{The Target Numbers - What You Need}
		\scriptsize
		\begin{block}{Publication Targets}
			\begin{itemize}
				\item \textbf{20-50 total publications} (journals + conferences + book chapters)
				\item \textbf{20-35 review papers} (high citation generators)
				\item \textbf{5-10 conference proceedings} (fast publication)
			\end{itemize}
		\end{block}
		\begin{block}{Citation Targets}
			\begin{itemize}
				\item \textbf{1000-2000 total Google Scholar citations}
				\item \textbf{1-10 .gov citations} (Federal Reserve, ERIC, Science.gov, BLS, NIH)
				\item \textbf{10-20 .edu citations} (Harvard, MIT, Stanford, Berkeley)
				\item \textbf{h-index: 15-25+}
			\end{itemize}
		\end{block}
	\end{frame}
	
	% ============================================================
	% SLIDE 3: WHY IMPACT > JOURNAL NAME
	% ============================================================
	\begin{frame}{Why Impact Matters More Than Journal Name}
		\scriptsize
		\begin{block}{The Truth About Journal Prestige}
			\begin{itemize}
				\item A paper in a mid-tier journal with 100 citations > paper in Nature with 0 citations
				\item USCIS cares about CITATIONS - proof your work is being used
				\item .gov citations from mid-tier journals = as valuable as top-tier
				\item Focus on getting your work READ and CITED, not where it appears
			\end{itemize}
		\end{block}
		\begin{block}{Where to Focus Your Energy - gov angle matters the most}
			\begin{itemize}
				\item \textbf{Content quality:} Solve real problems, address government priorities
				\item \textbf{Distribution:} Upload everywhere (arXiv, SSRN, OSF, ResearchGate)
				\item \textbf{Promotion:} Share on LinkedIn, email to researchers, present at conferences
				\item \textbf{Government alignment:} Topics that federal agencies care about
			\end{itemize}
		\end{block}
	\end{frame}
	
	% ============================================================
	% SLIDE 4: GOVERNMENT ANGLE - THE MOST IMPORTANT FACTOR
	% ============================================================
	\begin{frame}{Government Angle - The Most Important Factor}
		\scriptsize
		\begin{block}{Why Government Impact Matters Most}
			\begin{itemize}
				\item .gov citations are the SINGLE strongest evidence
				\item Federal agencies cite work that serves national interests
				\item One .gov citation can outweigh 50 regular citations
			\end{itemize}
		\end{block}
		\begin{block}{Federal Agencies to Target}
			\centering
			\begin{tabular}{|p{5cm}|p{5cm}|}
				\hline
				\textbf{Agency} & \textbf{Your Research Should Address} \\
				\hline
				Federal Reserve & Financial stability, payment systems, AI in banking \\
				\hline
				Department of Treasury & Crypto regulation, sanctions, fraud detection \\
				\hline
				Department of Labor & Workforce development, job training, AI displacement \\
				\hline
				Department of Education & AI literacy, digital skills, online learning \\
				\hline
				Department of Homeland Security & Cybersecurity, critical infrastructure \\
				\hline
				NIH/PubMed & Health AI, medical data, patient safety \\
				\hline
			\end{tabular}
		\end{block}
	\end{frame}
	
	% ============================================================
	% SLIDE 5: 50 PAPERS - THE VOLUME STRATEGY
	% ============================================================
	\begin{frame}{50 Papers - The High-Volume Strategy}
		\scriptsize
		\begin{block}{Why Volume Matters}
			\begin{itemize}
				\item More papers = more chances to be cited
				\item Each paper can generate 20-50 citations over time
				\item Volume demonstrates sustained research activity
				\item Manual submission to top indexing: MPRA, EconPaper, SSRN, ERIC, OSF, Sciexplorer
			\end{itemize}
		\end{block}
		\begin{block}{How to Achieve 50 Papers}
			\begin{itemize}
				\item \textbf{Write review papers:} 10-15 review papers (each cites 100+ refs, gets many citations)
				\item \textbf{Conference papers:} 15-20 conference proceedings (faster publication)
				\item \textbf{Collaborate:} Each collaboration adds 2-3 co-authored papers
				\item \textbf{Extend conference papers:} Conference $\rightarrow$ extended version $\rightarrow$ journal
				\item \textbf{Book chapters:} 5-8 invited chapters (editors reach out after you publish)
			\end{itemize}
		\end{block}
	\end{frame}
	
	% ============================================================
	% SLIDE 6: REVIEW PAPERS - YOUR CITATION ENGINE
	% ============================================================
	\begin{frame}{Review Papers - Your Citation Engine}
		\scriptsize
		\begin{block}{Why Review Papers Are Critical}
			\begin{itemize}
				\item Review papers get 2-5x MORE citations than original research
				\item Anyone starting research in the field WILL cite your review
				\item Stay relevant and citable for years
				\item Can be written without new experiments
			\end{itemize}
		\end{block}
		\begin{block}{How Many Review Papers You Need}
			\begin{itemize}
				\item Target: \textbf{10-15 review papers}
				\item Each review paper: 150-300 references cited
				\item Each review paper: Expected 100-300 citations
				\item 10 review papers × 150 citations = 1500 citations
			\end{itemize}
		\end{block}
		\begin{block}{Review Paper Topics (Government Focus)}
			\begin{itemize}
				\item "AI in Financial Stability: A Systematic Review" (Federal Reserve)
				\item "Cybersecurity for Critical Infrastructure" (DHS)
				\item "Digital Workforce Training: Methods and Outcomes" (DOL)
				\item "AI Literacy in K-12 Education" (Department of Education)
			\end{itemize}
		\end{block}
	\end{frame}
	
	% ============================================================
	% SLIDE 7: RAPID PUBLICATION VENUES
	% ============================================================
	\begin{frame}{Rapid Publication Venues - Build Volume Fast}
		\scriptsize
		\begin{block}{Fastest Publication Options}
			\begin{itemize}
				\item \textbf{MDPI journals:} Fast submission to publication
				\item \textbf{Elsevier open access:} Rapid turnaround
				\item \textbf{IEEE conferences:} Quick from submission to presentation
				\item \textbf{Springer conference proceedings:} Efficient process
			\end{itemize}
		\end{block}
		\begin{block}{Low-Cost Rapid Journals}
			\begin{itemize}
				\item MDPI: Some journals have low APCs or waivers
				\item Hindawi: Negotiate waivers available
				\item PeerJ: Tiered pricing options
				\item Directory of Open Access Journals: Filter by \$0 options
			\end{itemize}
		\end{block}
	\end{frame}
	
	% ============================================================
	% SLIDE 8: ZERO-COST PUBLICATION STRATEGY
	% ============================================================
	\begin{frame}{Zero-Cost Publication Strategy}
		\scriptsize
		\begin{block}{No-Cost Publication Options}
			\begin{itemize}
				\item \textbf{Traditional subscription journals:} Author pays nothing (IEEE Transactions, many Elsevier subscription journals)
				\item \textbf{University journals:} Many university presses charge no fee
				\item \textbf{Society journals:} IEEE, ACM, INFORMS member rates often \$0
				\item \textbf{Conference proceedings:} Registration only
				\item \textbf{Book chapters:} Usually no fee (invited)
				\item \textbf{Preprints:} Always free (arXiv, SSRN, OSF)
			\end{itemize}
		\end{block}
		\begin{block}{How to Find No-Cost Journals}
			\begin{enumerate}
				\item Go to Scopus Sources
				\item Search your field
				\item Check "Open Access" filter - select "All Open Access"
				\item Or select "Hybrid" journals (subscription track = no fee)
				\item Email editor to confirm no fee for subscription track
			\end{enumerate}
		\end{block}
	\end{frame}
	
	% ============================================================
	% SLIDE 9: CONFERENCE PROCEEDINGS - VOLUME BUILDER
	% ============================================================
	\begin{frame}{Conference Proceedings - High-Volume Strategy}
		\scriptsize
		\begin{block}{Why Conferences for Volume}
			\begin{itemize}
				\item Faster than journals
				\item Lower cost than many journals
				\item Can publish multiple conference papers
				\item Easy to extend into journal papers later
			\end{itemize}
		\end{block}
		\begin{block}{Conference Proceedings Indexed in Scopus}
			\begin{itemize}
				\item \textbf{IEEE Xplore:} All IEEE conferences (strongest)
				\item \textbf{Springer LNCS/LNAI/CCIS:} High-quality proceedings
				\item \textbf{ACM Digital Library:} ACM conferences
				\item \textbf{Elsevier Procedia:} Procedia Computer Science, Procedia Economics
				\item \textbf{CEUR Workshop Proceedings:} Free, Scopus indexed for some
			\end{itemize}
		\end{block}
		\begin{block}{Conference Strategy}
			\begin{itemize}
				\item Submit multiple conference papers
				\item Each conference = 1 publication, 5-20 eventual citations
				\item Present = networking for collaborations
			\end{itemize}
		\end{block}
	\end{frame}
	
	% ============================================================
	% SLIDE 10: COLLABORATION FOR VOLUME
	% ============================================================
	\begin{frame}{Collaboration - Force Multiplier for Volume}
		\scriptsize
		\begin{block}{Why Collaborate}
			\begin{itemize}
				\item Each collaboration adds multiple co-authored papers
				\item You don't have to write everything yourself
				\item Shared workload = more publications for everyone
			\end{itemize}
		\end{block}
		\begin{block}{How to Find Collaborators}
			\begin{itemize}
				\item \textbf{ResearchGate:} Message researchers in your field
				\item \textbf{LinkedIn:} Connect with academics and industry researchers
				\item \textbf{Conferences:} Meet potential collaborators in person
				\item \textbf{Twitter/X:} Follow and engage with researchers
				\item \textbf{Former colleagues:} Reach out to past coworkers
			\end{itemize}
		\end{block}
		\begin{block}{Collaboration Structure}
			\begin{itemize}
				\item You lead: Write first draft, handle submission (you get first author)
				\item Partner leads: They write, you contribute (you get co-author)
				\item Multiple small collaborations = many co-authored papers
			\end{itemize}
		\end{block}
	\end{frame}
	
	% ============================================================
	% SLIDE 11: PREPRINTS - CITATIONS BEFORE PUBLICATION
	% ============================================================
	\begin{frame}{Preprints - Citations Before Journal Publication}
		\scriptsize
		\begin{block}{Why Preprints Are Essential}
			\begin{itemize}
				\item Citations start the DAY you upload
				\item Google Scholar indexes preprints
				\item Preprints get DOI and become citable immediately
				\item Can be cited by government reports before journal publication
			\end{itemize}
		\end{block}
		\begin{block}{Where to Upload Every Paper}
			\begin{enumerate}
				\item \textbf{arXiv.org} (Very hard very reputed CS, physics, math, finance)
				\item \textbf{SSRN} (economics, finance, law)
				\item \textbf{OSF Preprints} (interdisciplinary, DOI from OSF)
				\item \textbf{ResearchGate} (upload full text)
			\end{enumerate}
		\end{block}
		\begin{block}{Preprint Citation Impact}
			\begin{itemize}
				\item Upload to arXiv (or MDPI preprints or SSRN)
				\item Paper indexed in Google Scholar
				\item First citation (someone reads your preprint)
				\item Journal publication (citations already started)
				\item Many citations from preprint version
			\end{itemize}
		\end{block}
	\end{frame}
	
	% ============================================================
	% SLIDE 12: .GOV CITATIONS - FEDERAL RESERVE
	% ============================================================
	\begin{frame}{.gov Citations - Federal Reserve (FEDS Series)}
		\scriptsize
		\begin{block}{Targeting Federal Reserve Citations}
			\begin{itemize}
				\item Federal Reserve economists read SSRN and arXiv preprints
				\item Your preprint can be cited BEFORE journal publication
				\item FEDS papers have DOIs and are .gov domain
			\end{itemize}
		\end{block}
		\begin{block}{Topics That Get Fed Citations}
			\begin{itemize}
				\item Machine learning for financial stability monitoring
				\item Payment systems and FedNow adoption
				\item Climate risk and financial institutions
				\item AI in banking supervision and compliance
				\item Cryptocurrency and digital asset regulation
				\item Stress testing methodologies
			\end{itemize}
		\end{block}
		\begin{block}{How to Maximize Fed Citation Chance}
			\begin{enumerate}
				\item Upload to SSRN (Fed economists use SSRN heavily)
				\item Include "Federal Reserve" or "financial stability" in keywords
				\item Cite Fed papers in your work (they notice)
				\item Email your SSRN link to Fed economists
			\end{enumerate}
		\end{block}
	\end{frame}
	
	% ============================================================
	% SLIDE 13: .GOV CITATIONS - ERIC (DEPARTMENT OF EDUCATION)
	% ============================================================
	\begin{frame}{.gov Citations - ERIC (Department of Education)}
		\scriptsize
		\begin{block}{What is ERIC?}
			\begin{itemize}
				\item Education Resources Information Center
				\item .gov domain of U.S. Department of Education
				\item Indexes education, workforce, training research
			\end{itemize}
		\end{block}
		\begin{block}{How Papers Get into ERIC}
			\begin{enumerate}
				\item Publish in ERIC-indexed journal
				\item Journal automatically sends content to ERIC
				\item Your paper appears on ERIC with .gov URL
				\item Your work is discoverable through .gov portal
			\end{enumerate}
		\end{block}
		\begin{block}{Topics for ERIC Indexing}
			\begin{itemize}
				\item AI literacy and digital skills training
				\item Workforce development and retraining programs
				\item Online learning effectiveness research
				\item Veteran education and job training
				\item STEM education and diversity
				\item Community college workforce programs
			\end{itemize}
		\end{block}
	\end{frame}
	
	% ============================================================
	% SLIDE 14: .GOV CITATIONS - SCIENCE.GOV
	% ============================================================
	\begin{frame}{.gov Citations - Science.gov}
		\scriptsize
		\begin{block}{What is Science.gov?}
			\begin{itemize}
				\item Official U.S. government search portal for science
				\item Searches 60+ federal databases
				\item Managed by Department of Energy
			\end{itemize}
		\end{block}
		\begin{block}{Pathways to Science.gov}
			\begin{itemize}
				\item \textbf{PubMed} (NIH) $\rightarrow$ Science.gov
				\item \textbf{ERIC} (Dept of Education) $\rightarrow$ Science.gov
				\item \textbf{NASA ADS} $\rightarrow$ Science.gov
				\item \textbf{DOE OSTI} $\rightarrow$ Science.gov
				\item \textbf{USDA} $\rightarrow$ Science.gov
			\end{itemize}
		\end{block}
		\begin{block}{Strategy for Science.gov Indexing}
			\begin{itemize}
				\item Publish in PubMed-indexed journals for health/AI topics
				\item Publish in ERIC-indexed journals for education/workforce
				\item Upload to DOE OSTI if your research has energy relevance
				\item Once indexed, your work is searchable on .gov domain
			\end{itemize}
		\end{block}
	\end{frame}
	
	% ============================================================
	% SLIDE 15: .GOV CITATIONS - BUREAU OF LABOR STATISTICS
	% ============================================================
	\begin{frame}{.gov Citations - Bureau of Labor Statistics}
		\scriptsize
		\begin{block}{Targeting BLS Citations}
			\begin{itemize}
				\item BLS publishes Monthly Labor Review (MLR)
				\item Also publishes working papers and reports
				\item Citations from BLS = .gov domain evidence
			\end{itemize}
		\end{block}
		\begin{block}{Topics for BLS Citations}
			\begin{itemize}
				\item Employment projections and labor market trends
				\item Impact of AI on employment and wages
				\item Workforce demographics and participation rates
				\item Occupational outlook and skills demand
				\item Gig economy and alternative work arrangements
			\end{itemize}
		\end{block}
		\begin{block}{How to Get Cited by BLS}
			\begin{enumerate}
				\item Publish on labor economics topics
				\item Upload to SSRN (BLS researchers use SSRN)
				\item Submit comments to BLS RFIs on Regulations.gov
				\item Cite BLS data in your papers (they notice)
			\end{enumerate}
		\end{block}
	\end{frame}
	
	% ============================================================
	% SLIDE 16: .GOV CITATIONS - NIH/PUBMED
	% ============================================================
	\begin{frame}{.gov Citations - NIH / PubMed}
		\scriptsize
		\begin{block}{Targeting NIH Citations}
			\begin{itemize}
				\item PubMed is .gov domain (nih.gov)
				\item Citations from NIH-funded research are .gov
				\item Health-related AI research is high priority
			\end{itemize}
		\end{block}
		\begin{block}{Topics for PubMed Citations}
			\begin{itemize}
				\item AI for clinical decision support
				\item Machine learning for patient safety
				\item Digital health and telemedicine
				\item Health data privacy and security
				\item Medical image analysis
				\item Public health surveillance
			\end{itemize}
		\end{block}
		\begin{block}{How to Get PubMed Citations}
			\begin{enumerate}
				\item Publish in PubMed-indexed journals
				\item Upload to medRxiv preprint server
				\item Collaborate with NIH-funded researchers
				\item Focus on health applications of your methods
			\end{enumerate}
		\end{block}
	\end{frame}
	
	% ============================================================
	% SLIDE 17: .EDU CITATIONS - TOP UNIVERSITIES
	% ============================================================
	\begin{frame}{.edu Citations - Top US Universities}
		\scriptsize
		\begin{block}{Why .edu Citations Are Valuable}
			\begin{itemize}
				\item Verifiable .edu domain
				\item Demonstrates academic recognition in the U.S.
				\item Harvard, MIT, Stanford citations are strongest
			\end{itemize}
		\end{block}
		\begin{block}{How to Get .edu Citations}
			\begin{enumerate}
				\item Publish in reputable journals (Scopus/WoS)
				\item Upload to institutional repositories (Harvard DASH, MIT DSpace)
				\item Present at university seminars (virtual or in-person)
				\item Collaborate with US university researchers
				\item Create course materials that professors adopt
			\end{enumerate}
		\end{block}
		\begin{block}{Tracking .edu Citations}
			\begin{itemize}
				\item Search Google Scholar with site:.edu "your paper title"
				\item Save screenshots of results with .edu URLs
				\item Document citations from specific universities
			\end{itemize}
		\end{block}
	\end{frame}
	
	% ============================================================
	% SLIDE 18: EBSCO HOST - MAJOR COMMERCIAL DATABASE
	% ============================================================
	\begin{frame}{EBSCO Host - Massive Academic Reach}
		\scriptsize
		\begin{block}{What is EBSCO?}
			\begin{itemize}
				\item Largest academic database aggregator worldwide
				\item Serves thousands of universities globally
				\item Indexes thousands of journals across all disciplines
				\item Content appears in university library search results
			\end{itemize}
		\end{block}
		\begin{block}{Why EBSCO Citations Matter}
			\begin{itemize}
				\item Your paper appears when students \& researchers search their university library
				\item More discoverability = more citations
				\item Many Scopus-indexed journals are ALSO in EBSCO
				\item EBSCO discovery services power university search portals
			\end{itemize}
		\end{block}
		\begin{block}{How to Get Indexed in EBSCO}
			\begin{itemize}
				\item Publish in journals that EBSCO indexes (check Journal Finder)
				\item Many low-cost journals are EBSCO-indexed
				\item EBSCO indexes conference proceedings and book chapters too
				\item Ask your journal editor: "Is this journal in EBSCO?"
			\end{itemize}
		\end{block}
	\end{frame}
	
	% ============================================================
	% SLIDE 19: PROQUEST - DISSERTATIONS AND BEYOND
	% ============================================================
	\begin{frame}{ProQuest - Beyond Traditional Journals}
		\scriptsize
		\begin{block}{What is ProQuest?}
			\begin{itemize}
				\item Major academic database alongside EBSCO
				\item Indexes dissertations, journals, newspapers, conference papers
				\item Available in most US university libraries
				\item ProQuest Central is most widely used
			\end{itemize}
		\end{block}
		\begin{block}{ProQuest Citation Benefits}
			\begin{itemize}
				\item Citations from .edu domains via university library access
				\item Dissertations citing your work = .edu citations
				\item Conference proceedings indexed in ProQuest get more visibility
				\item Government reports indexed here too
			\end{itemize}
		\end{block}
		\begin{block}{How to Get into ProQuest}
			\begin{itemize}
				\item Conference proceedings often automatically go to ProQuest
				\item Dissertations that cite you are in PQDT (ProQuest Dissertations)
				\item Publish in ProQuest-indexed journals
				\item Upload to institutional repositories (university libraries submit)
			\end{itemize}
		\end{block}
	\end{frame}
	
	% ============================================================
	% SLIDE 20: PEER REVIEW - BUILDING RECOGNITION
	% ============================================================
	\begin{frame}{Peer Review - Building Recognition}
		\scriptsize
		\begin{block}{Target: Multiple Peer Reviews}
			\begin{itemize}
				\item Start with a few reviews
				\item Grow to become established reviewer
				\item Progress to senior reviewer status
			\end{itemize}
		\end{block}
		\begin{block}{How to Get Invited}
			\begin{enumerate}
				\item Publish papers in a journal (editors notice)
				\item Register as reviewer on ScholarOne and Editorial Manager
				\item Email journal editors: "I have published in your journal and am available to review"
				\item Accept invitations when possible
			\end{enumerate}
		\end{block}
		\begin{block}{Review Documentation}
			\begin{itemize}
				\item Save all invitation emails
				\item Save completion confirmations
				\item Web of Science reviewer profile (automatic tracking)
				\item ORCID reviewer recognition
			\end{itemize}
		\end{block}
	\end{frame}
	
	% ============================================================
	% SLIDE 21: AWARDS - BUILDING RECOGNITION
	% ============================================================
	\begin{frame}{Awards - Building Recognition}
		\scriptsize
		\begin{block}{Award Opportunities}
			\begin{itemize}
				\item \textbf{Best Paper Awards} at conferences you attend
				\item \textbf{Young Researcher Awards} from professional societies
				\item \textbf{Travel Grants} (award for attending)
				\item \textbf{Outstanding Reviewer Awards} from journals
				\item \textbf{Teaching Awards} from your employer
			\end{itemize}
		\end{block}
		\begin{block}{Awards to Target}
			\begin{itemize}
				\item IEEE Outstanding Young Professional
				\item ACM Student Research Competition
				\item Conference Best Paper Award
				\item Journal Outstanding Reviewer Award
			\end{itemize}
		\end{block}
	\end{frame}
	
	% ============================================================
	% SLIDE 22: BOOK CHAPTERS - ADDING VOLUME
	% ============================================================
	\begin{frame}{Book Chapters - Adding Publication Volume}
		\scriptsize
		\begin{block}{How to Get Book Chapter Invitations}
			\begin{itemize}
				\item Publish journal papers first
				\item Editors of edited volumes will find you on Google Scholar
				\item They email invitation to contribute a chapter
				\item Accept invitations when possible
			\end{itemize}
		\end{block}
		\begin{block}{Book Chapter Benefits}
			\begin{itemize}
				\item Most book chapters have no publication fee
				\item Target multiple chapters
				\item Publishers include Springer, Elsevier, Wiley
			\end{itemize}
		\end{block}
	\end{frame}
	
	% ============================================================
	% SLIDE 23: PREPRINT SERVERS - EVERY PAPER GOES HERE
	% ============================================================
	\begin{frame}{Preprint Servers - Every Paper Must Go Here}
		\scriptsize
		\begin{block}{Upload Every Paper to These Servers}
			\begin{enumerate}
				\item \textbf{arXiv.org} (Very hard, CS, physics, math, finance, statistics)
				\item \textbf{SSRN} (for economics, finance, law, social sciences)
				\item \textbf{OSF Preprints} (for interdisciplinary, gets DOI)
				\item \textbf{ResearchGate} (full text upload)
			\end{enumerate}
		\end{block}
		\begin{block}{Why This Matters for Citations}
			\begin{itemize}
				\item Each server = additional discovery pathway
				\item Each server = potential new reader = potential citation
				\item Multiple servers × multiple papers = many distribution points
				\item Citations start the day you upload
			\end{itemize}
		\end{block}
		\begin{block}{Preprint Citation Evidence for Portfolio}
			\begin{itemize}
				\item Screenshot of arXiv/SSRN download counts
				\item Citation of your preprint in other papers
				\item Google Scholar shows preprint citations
			\end{itemize}
		\end{block}
	\end{frame}
	
	% ============================================================
	% SLIDE 24: LINKEDIN FOR RESEARCH VISIBILITY
	% ============================================================
	\begin{frame}{LinkedIn - Research Promotion}
		\scriptsize
		\begin{block}{LinkedIn Strategy for Citations}
			\begin{itemize}
				\item Every time a paper is published/posted, share on LinkedIn
				\item Tag co-authors and relevant researchers
				\item Use hashtags: \#AI \#MachineLearning \#Cybersecurity \#FinTech
				\item Engage with comments (algorithm boosts visibility)
			\end{itemize}
		\end{block}
		\begin{block}{Results from Consistent LinkedIn Posting}
			\begin{itemize}
				\item More downloads when shared on LinkedIn
				\item More downloads = more citations
				\item Builds network of researchers who will cite you
			\end{itemize}
		\end{block}
	\end{frame}
	
	% ============================================================
	% SLIDE 25: TWITTER/X FOR RESEARCH NETWORKING
	% ============================================================
	\begin{frame}{Twitter/X - Research Networking}
		\scriptsize
		\begin{block}{Why Twitter/X for Researchers}
			\begin{itemize}
				\item Active academic community (hashtags: \#AcademicTwitter, \#AI, \#MachineLearning)
				\item Journal editors and conference organizers are on Twitter
				\item Preprints shared here get more citations
			\end{itemize}
		\end{block}
		\begin{block}{Twitter Strategy}
			\begin{itemize}
				\item Share every paper with 1-2 sentence summary
				\item Use relevant hashtags
				\item Tag journal/conference account if applicable
				\item Retweet and engage with other researchers' papers
			\end{itemize}
		\end{block}
	\end{frame}
	
	% ============================================================
	% SLIDE 26: ZOTERO FOR LITERATURE REVIEW - SCALE UP
	% ============================================================
	\begin{frame}{Zotero - Literature Review at Scale}
		\scriptsize
		\begin{block}{Why Zotero for 50 Papers}
			\begin{itemize}
				\item Each review paper cites 100-300 references
				\item 10-15 review papers = many references managed
				\item Zotero organizes everything automatically
			\end{itemize}
		\end{block}
		\begin{block}{Zotero Setup for Volume}
			\begin{itemize}
				\item Create collections for each paper
				\item Use tags for organization
				\item Install Better BibTeX for LaTeX export
				\item Use browser connector to save papers instantly
			\end{itemize}
		\end{block}
		\begin{block}{Literature Review Workflow}
			\begin{enumerate}
				\item Search Google Scholar for relevant papers
				\item Save to Zotero with one click
				\item Export .bib file for your LaTeX paper
				\item Write review paper citing these references
			\end{enumerate}
		\end{block}
	\end{frame}
	
	% ============================================================
	% SLIDE 27: LATEX FOR 50 PAPERS - TEMPLATE SYSTEM
	% ============================================================
	\begin{frame}{LaTeX - Template System for 50 Papers}
		\scriptsize
		\begin{block}{Why LaTeX for Volume}
			\begin{itemize}
				\item Create templates - reuse for every paper
				\item Automatic bibliography (BibTeX) - no manual formatting
				\item Consistent formatting across many papers
				\item Free (Overleaf or TeXstudio)
			\end{itemize}
		\end{block}
		\begin{block}{Template Structure}
			\centering
			\begin{tabular}{ll}
				\multicolumn{2}{l}{\texttt{paper\_template.tex}} \\
				\texttt{$\vdash$ preamble.tex} & (packages, commands) \\
				\texttt{$\vdash$ title.tex} & (title, authors, abstract) \\
				\texttt{$\vdash$ introduction.tex} & \\
				\texttt{$\vdash$ literature\_review.tex} & \\
				\texttt{$\vdash$ methodology.tex} & \\
				\texttt{$\vdash$ results.tex} & \\
				\texttt{$\vdash$ discussion.tex} & \\
				\texttt{$\vdash$ conclusion.tex} & \\
				\texttt{$\vdash$ references.bib} & (from Zotero) \\
			\end{tabular}
		\end{block}
		\begin{itemize}
			\item Copy template folder for each new paper
			\item Modify content, keep structure
			\item Saves time per paper
		\end{itemize}
	\end{frame}
	
	% ============================================================
	% SLIDE 28: AI TOOLS FOR DRAFTING - ACCELERATE WRITING
	% ============================================================
	\begin{frame}{AI Tools - Accelerate Writing to 50 Papers}
		\scriptsize
		\begin{block}{Using AI for Volume}
			\begin{itemize}
				\item \textbf{Literature review summaries:} AI summarizes many papers quickly
				\item \textbf{First draft generation:} Provide outline, AI writes draft
				\item \textbf{Editing and polishing:} Improve clarity and grammar
				\item \textbf{Reference formatting:} Ensure BibTeX is correct
			\end{itemize}
		\end{block}
		\begin{block}{Sample AI Prompt for Review Paper}
			\begin{center}
				\ttfamily
				Write a literature review on \\
				"Machine Learning for Financial Stability" \\
				including sections on: (1) anomaly detection, \\
				(2) stress testing, (3) payment systems. \\
				Cite numerous references.
			\end{center}
		\end{block}
	\end{frame}
	
	% ============================================================
	% SLIDE 29: GITHUB FOR COLLABORATION
	% ============================================================
	\begin{frame}{GitHub - Collaboration for Volume}
		\scriptsize
		\begin{block}{Why GitHub for Research}
			\begin{itemize}
				\item Version control for all papers
				\item Collaborate with co-authors without emailing files
				\item Public profile shows research activity
				\item Code and data can be shared
			\end{itemize}
		\end{block}
		\begin{block}{GitHub Workflow}
			\begin{enumerate}
				\item Create repository for each paper
				\item Push LaTeX code, figures, .bib file
				\item Add co-authors as collaborators
				\item Track changes (who wrote what)
				\item Merge contributions automatically
			\end{enumerate}
		\end{block}
	\end{frame}
	
	% ============================================================
	% SLIDE 30: 50 PAPERS - REALISTIC BREAKDOWN BY TYPE
	% ============================================================
	\begin{frame}{50 Papers - Realistic Breakdown by Type}
		\scriptsize
		\begin{block}{Target Mix for 50 Papers}
			\centering
			\begin{tabular}{|p{4cm}|p{3cm}|p{3cm}|}
				\hline
				\textbf{Paper Type} & \textbf{Number} & \textbf{Citations Each} \\
				\hline
				Review Papers & 12 & 150-300 \\
				\hline
				Original Research (Journals) & 15 & 30-80 \\
				\hline
				Conference Proceedings & 15 & 10-30 \\
				\hline
				Book Chapters & 5 & 10-30 \\
				\hline
				Communications & 3 & 5-20 \\
				\hline
				\textbf{Total} & \textbf{50} & \textbf{1000-2000+} \\
				\hline
			\end{tabular}
		\end{block}
	\end{frame}
	
	% ============================================================
	% SLIDE 31: SAMPLE PAPER TOPICS - FINANCIAL STABILITY
	% ============================================================
	\begin{frame}{Sample Paper Topics - Financial Stability}
		\scriptsize
		\begin{block}{Review Papers}
			\begin{enumerate}
				\item "Machine Learning for Financial Stability: A Systematic Review"
				\item "AI in Banking Supervision: Methods and Applications"
				\item "Cryptocurrency Regulation: A Review of Global Approaches"
				\item "Stress Testing Methodologies: Evolution and Future Directions"
				\item "Payment Systems Security: Threats and Countermeasures"
			\end{enumerate}
		\end{block}
		\begin{block}{Original Research Papers}
			\begin{enumerate}
				\item "Detecting Anomalies in Real-Time Payment Systems Using Graph Neural Networks"
				\item "Forecasting Bank Runs with Social Media Sentiment Analysis"
				\item "Explainable AI for Anti-Money Laundering Compliance"
				\item "Blockchain for Cross-Border Payments: Efficiency Analysis"
				\item "Machine Learning for Credit Risk Assessment in Community Banks"
			\end{enumerate}
		\end{block}
	\end{frame}
	
	% ============================================================
	% SLIDE 32: SAMPLE PAPER TOPICS - WORKFORCE DEVELOPMENT
	% ============================================================
	\begin{frame}{Sample Paper Topics - Workforce Development}
		\scriptsize
		\begin{block}{Review Papers}
			\begin{enumerate}
				\item "AI and Workforce Displacement: A Systematic Review of Retraining Programs"
				\item "Digital Literacy Frameworks for Adult Learners"
				\item "Online Learning Effectiveness for Workforce Training"
				\item "Veteran Transition to Tech Careers: Best Practices"
				\item "Community College Workforce Programs: Models and Outcomes"
			\end{enumerate}
		\end{block}
		\begin{block}{Original Research Papers}
			\begin{enumerate}
				\item "Evaluating AI Literacy Curriculum for Community College Students"
				\item "Predicting Workforce Training Outcomes Using Machine Learning"
				\item "Gamification for Digital Skills Training: Randomized Controlled Trial"
				\item "Employer Perspectives on AI Credentials and Certifications"
				\item "Economic Impact of Digital Upskilling Programs on Local Communities"
			\end{enumerate}
		\end{block}
	\end{frame}
	
	% ============================================================
	% SLIDE 33: SAMPLE PAPER TOPICS - CYBERSECURITY
	% ============================================================
	\begin{frame}{Sample Paper Topics - Cybersecurity}
		\scriptsize
		\begin{block}{Review Papers}
			\begin{enumerate}
				\item "AI for Cybersecurity: A Systematic Review of Detection Methods"
				\item "Critical Infrastructure Protection: Frameworks and Gaps"
				\item "Zero-Trust Architecture: Implementation Challenges and Solutions"
				\item "Ransomware Trends and Defense Strategies"
				\item "Supply Chain Cybersecurity: Risk Assessment Methodologies"
			\end{enumerate}
		\end{block}
		\begin{block}{Original Research Papers}
			\begin{enumerate}
				\item "Detecting Phishing Attacks Using Large Language Models"
				\item "Anomaly Detection in Industrial Control Systems with Autoencoders"
				\item "Federated Learning for Collaborative Threat Intelligence"
				\item "Automated Vulnerability Assessment Using Graph Neural Networks"
				\item "Privacy-Preserving Machine Learning for Healthcare Data"
			\end{enumerate}
		\end{block}
	\end{frame}
	
	% ============================================================
	% SLIDE 34: SAMPLE PAPER TOPICS - AI IN GOVERNMENT
	% ============================================================
	\begin{frame}{Sample Paper Topics - AI in Government}
		\scriptsize
		\begin{block}{Review Papers}
			\begin{enumerate}
				\item "AI Governance Frameworks: A Comparative Analysis of Federal Agencies"
				\item "Algorithmic Bias Assessment Methods for Government Systems"
				\item "Explainable AI for Regulatory Compliance"
				\item "AI Procurement Best Practices for Public Sector"
				\item "Privacy-Preserving AI for Government Data Sharing"
			\end{enumerate}
		\end{block}
		\begin{block}{Original Research Papers}
			\begin{enumerate}
				\item "Automating FOIA Redaction Using NLP"
				\item "AI for Regulatory Comment Analysis at Scale"
				\item "Predictive Models for Government Program Outcomes"
				\item "Chatbots for Citizen Service Delivery: Effectiveness Study"
				\item "AI-Assisted Policy Impact Assessment: Case Studies"
			\end{enumerate}
		\end{block}
	\end{frame}
	
	% ============================================================
	% SLIDE 35: MAXIMIZING .GOV CITATIONS - ACTION PLAN
	% ============================================================
	\begin{frame}{Maximizing .gov Citations - Action Plan}
		\scriptsize
		\begin{block}{Regular Actions for .gov Citations}
			\begin{itemize}
				\item Check Regulations.gov for new RFIs in your area
				\item Submit comments to RFIs citing your own research
				\item Search for new .gov citations to your papers
				\item Email your SSRN/arXiv links to agency researchers
			\end{itemize}
		\end{block}
		\begin{block}{Agency Outreach Strategy}
			\begin{enumerate}
				\item Identify researchers at Federal Reserve, BLS, ERIC, NIH
				\item Find their email (agency websites, Google Scholar)
				\item Send professional email sharing your relevant preprint
				\item No requests, just sharing - they may cite if useful
			\end{enumerate}
		\end{block}
	\end{frame}
	
	% ============================================================
	% SLIDE 36: CITATION TRACKING DASHBOARD
	% ============================================================
	\begin{frame}{Citation Tracking Dashboard}
		\scriptsize
		\begin{block}{What to Track Regularly}
			\centering
			\begin{tabular}{|p{4cm}|p{5.5cm}|}
				\hline
				\textbf{Metric} & \textbf{How to Track} \\
				\hline
				Total Google Scholar citations & Google Scholar profile \\
				\hline
				New citations & Google Scholar "Cited by" alerts \\
				\hline
				.gov citations & Google Scholar search: site:.gov "your name" \\
				\hline
				.edu citations & Google Scholar search: site:.edu "your name" \\
				\hline
				h-index & Google Scholar profile \\
				\hline
				Peer reviews completed & Web of Science / ORCID \\
				\hline
				Preprint downloads & arXiv/SSRN/OSF stats \\
				\hline
			\end{tabular}
		\end{block}
		\begin{block}{Sample Tracking Sheet}
			\centering
			\begin{tabular}{|c|c|c|c|c|}
				\hline
				\textbf{Papers} & \textbf{Citations} & \textbf{.gov} & \textbf{.edu} & \textbf{Reviews} \\
				\hline
				8 & 25 & 0 & 2 & 3 \\
				\hline
				15 & 80 & 1 & 5 & 8 \\
				\hline
				25 & 300 & 3 & 12 & 15 \\
				\hline
				35 & 800 & 5 & 25 & 25 \\
				\hline
				45 & 1500 & 7 & 40 & 35 \\
				\hline
				50 & 2000 & 8 & 50 & 50 \\
				\hline
			\end{tabular}
		\end{block}
	\end{frame}
	
	% ============================================================
	% SLIDE 37: COMMON MISTAKES - HIGH-VOLUME EDITION
	% ============================================================
	\begin{frame}{Common Mistakes - High-Volume Edition}
		\scriptsize
		\begin{block}{Mistakes That Kill Volume Strategy}
			\begin{itemize}
				\item \textcolor{red}{Only targeting top-tier journals} (rejections waste time)
				\item \textcolor{red}{Writing every paper alone} (collaboration = more papers)
				\item \textcolor{red}{Not using preprints} (delays citations)
				\item \textcolor{red}{Not tracking citations systematically} (can't prove impact)
				\item \textcolor{red}{Ignoring government alignment} (misses .gov citations)
			\end{itemize}
		\end{block}
		\begin{block}{Mistakes to Avoid - Peer Review}
			\begin{itemize}
				\item \textcolor{red}{Declining all review invitations} (misses recognition)
				\item \textcolor{red}{Not registering on journal portals} (no invitations)
			\end{itemize}
		\end{block}
	\end{frame}
	
	% ============================================================
	% SLIDE 38: FINAL CHECKLIST - 50 PAPERS / 2000 CITATIONS
	% ============================================================
	\begin{frame}{Final Checklist - 50 Papers | 2000 Citations Ready}
		\scriptsize
		\begin{block}{Publications Checklist}
			\begin{itemize}
				\item $\square$ 50 total publications (journals + conferences + chapters)
				\item $\square$ 12-15 review papers (high citation generators)
				\item $\square$ All papers uploaded to preprint servers
			\end{itemize}
		\end{block}
		\begin{block}{Citations Checklist}
			\begin{itemize}
				\item $\square$ 1000-2000+ Google Scholar citations
				\item $\square$ .gov citations documented
				\item $\square$ .edu citations documented
				\item $\square$ Citation growth trend (upward)
			\end{itemize}
		\end{block}
		\begin{block}{Recognition Checklist}
			\begin{itemize}
				\item $\square$ Many peer reviews completed
				\item $\square$ Web of Science reviewer profile
				\item $\square$ Awards (conference best paper, outstanding reviewer)
				\item $\square$ Book chapters
			\end{itemize}
		\end{block}
		\begin{block}{Profile Management Checklist}
			\begin{itemize}
				\item $\square$ Public Google Scholar profile with all papers
				\item $\square$ ORCID with publications and reviews
				\item $\square$ ResearchGate with full texts
				\item $\square$ LinkedIn with publications section
			\end{itemize}
		\end{block}
	\end{frame}
	
	% ============================================================
	% SLIDE 39: CONCLUSION - YOU CAN DO THIS
	% ============================================================
	\begin{frame}{Conclusion - You Can Build This Profile}
		\scriptsize
		\begin{block}{The Formula for 50 Papers \& 2000 Citations}
			\begin{enumerate}
				\item \textbf{Write review papers} (12-15) - each generates many citations
				\item \textbf{Collaborate widely} - each collaborator adds multiple papers
				\item \textbf{Use preprints} - citations start immediately
				\item \textbf{Target government priorities} - .gov citations are gold
				\item \textbf{Low-cost journals} - stay within budget
				\item \textbf{AI-assisted writing} - accelerate to 50 papers
				\item \textbf{Consistent effort} - consistency beats intensity
			\end{enumerate}
		\end{block}
		\begin{block}{The Path Forward}
			\begin{itemize}
				\item Build foundation with initial papers
				\item Momentum builds with more papers
				\item Exponential growth in citations
			\end{itemize}
		\end{block}
		\begin{block}{Final Word}
			\textcolor{blue}{Impact > Journal Name. Government Angle > Everything.}
			
			\textcolor{blue}{Volume creates citations. Citations create recognition.}
			
			\textcolor{green}{Start today. One paper at a time. 50 papers.}
		\end{block}
	\end{frame}
	
\end{document}